\section{
  \texorpdfstring{总结与展望} % 正文和目录中显示的标题
  {\thesection 总结与展望} % 书签页中显示的标题
}
\label{sec:总结与展望}

\subsection{
  \texorpdfstring{总结} % 正文和目录中显示的标题
  {\thesubsection 总结} % 书签页中显示的标题
}
\label{ssec:总结}

海洋工程中海上浮式平台在恶劣海况下极易发生波浪砰击、上浪等现象，对人员设备和平台结构可能造成严重危害。%
对于这一类强非线性波浪浮体相互作用问题，CFD数值模拟为了精确捕捉自由液面需要耗费巨大的计算资源。%
本文对势粘流耦合数值方法展开了探索性研究，提出了基于势流方法和无网格SPH方法的势粘流耦合模型。%
文中基于HOS方法、SPH方法和BEM方法，在二维上分别实现了HOS-SPH势粘流单向耦合方法、%
SPH-BEM势粘流双向耦合方法和HOS-SPH-BEM势粘流多模型耦合方法，%
并通过和其他数值方法以及文献中的试验结果对比验证了几类耦合方法的有效性和准确性。%
全文主要结论如下：%

(1)	本文提出了基于势流方法和无网格粒子法的耦合模型。%
该模型通过基于粒子间相互作用的边界条件，实现了势流方法向无网格粒子法速度、压力等流场信息的传递，%
同时基于流量监测网格实现了势粘流耦合过程中拉格朗日粒子法的流入/流出，%
从而实现了流场质量、动量的传递，并和其他数值方法以及试验结果对比验证了该耦合模型的有效性。%

(2)	基于所提出的势粘流耦合模型，本文实现了HOS-SPH单向耦合方法。%
该方法引入了势流HOS方法用于模拟波浪，并且采用耦合模型有效地将势流部分流场信息传递至了粘流部分，从而有效地提高波浪模拟质量和计算效率。%
通过规则波和不规则波模拟、规则波浮体上浪和聚焦波冲击浮体等算例，验证了HOS-SPH单向耦合方法对于计算效率的提高，%
并初步验证了HOS-SPH单向耦合方法对于强非线性波浪结构物相互作用问题的有效性。%
聚焦波冲击浮体算例中耦合方法的数值结果较好地实现了对试验中上浪过程的模拟和复现，并且比较准确地捕捉到了上浪过程中甲板所受到的冲击压力峰值。%

(3)	本文对HOS-SPH单向耦合方法中耦合区域宽度和粘流域宽度这两个重要参数进行了探索，并给出了一些经验参数。%
通过数值波浪模拟算例对不同耦合区域宽度下的结果进行分析，得出耦合区域宽度应选取为$D/\Delta x = 10 \sim 20$，当粒子分辨率较高时该宽度可适当增大。%
通过聚焦波冲击浮体算例对不同粘流域宽度下的结果进行分析，得出在该算例中浮体前端到耦合边界的合适距离约为1.5倍波长，%
而后端到耦合边界的距离约为2倍波长，从而能够保证所关注计算结果的准确性的情况下尽量减少计算成本。%
但由于不同问题中波浪和结构物的尺度、流固耦合的剧烈程度等都有所不同，%
粘流域宽度的选取应视研究问题而定，以保证强非线性现象和粘性效应显著的区域包含在粘流域内。%

(4)	针对单向耦合中耦合边界处反射波的问题，本文实现了SPH-BEM双向耦合方法。%
通过规则波和不规则波模拟，进行了SPH-BEM双向耦合方法的粒子分辨率收敛性分析，并讨论了势流域边界与耦合边界间的距离对波高结果的影响。%
结果表明当势流域边界与耦合边界的距离为2倍波长时波高结果基本收敛。%
波浪数值模拟验证了该方法能够有效地将波浪从SPH粘流域传递到BEM势流域而不在耦合边界上发生反射，从而减小SPH计算域的大小并提高计算效率。%

(5)	在HOS-SPH单向耦合方法和SPH-BEM双向耦合方法的基础上，本文提出了HOS-SPH-BEM多模型耦合方法。%
该方法将势流部分的波浪分解为入射波和辐射绕射波，并分别采用HOS方法和BEM方法求解，%
在保留HOS方法能够高效准确模拟非线性波浪的优势的同时实现了势粘流间的双向耦合。%
在聚焦波冲击浮体算例中，通过与HOS-SPH单向耦合方法结果以及其他结果进行对比，%
验证了HOS-SPH-BEM多模型耦合方法能够减小耦合边界反射波对计算结果的影响，从而进一步提高计算精度。

\subsection{
  \texorpdfstring{创新性工作} % 正文和目录中显示的标题
  {\thesubsection 创新性工作} % 书签页中显示的标题
}
\label{ssec:创新性工作}

本文以恶劣海况下的海上浮式平台为研究背景，针对强非线性波浪浮体相互作用问题中CFD方法难以快速准确计算的难题，%
对势粘流耦合数值方法展开了探索性研究，主要有以下创新性工作：%

(1)	针对波浪与浮体相互作用的数值模拟，提出了基于势流方法与无网格粒子法的耦合模型，初步探索了粘势流场间信息相互传递的难点。%

(2)	自主开发了二维的HOS-SPH单向耦合和HOS-SPH-BEM多模型耦合程序，实现了波浪与浮体相互作用的高效准确模拟。

\subsection{
  \texorpdfstring{展望} % 正文和目录中显示的标题
  {\thesubsection 展望} % 书签页中显示的标题
}
\label{ssec:展望}

本文对波浪和浮体结构物相互作用问题，进行了势粘流耦合方法的初步探索，但仍有大量工作有待开展，以下是作者对此方向的一些展望：%

(1)	本文所实现的HOS-SPH-BEM多模型双向耦合方法中所采用的边界元方法并非全非线性BEM，%
且未分析势粘流耦合过程中的动量、能量机制，因此在后续工作中需要对双向耦合方法进行深入的系统性分析和机理性研究。%

(2)	本文仅在二维上实现了所提出的势粘流耦合方法，且程序仅采用CPU并行运算，暂时难以应用于工程实际中。%
但耦合计算中的边界条件处理方式可以直接推广至三维，并且SPH方法十分适合于GPU并行运算，%
因此未来工作中可将该耦合方法发展至三维并采取GPU并行运算，以适应工程应用中对砰击、上浪等问题的研究。%
此外可以考虑将该耦合方法与水弹性方法等相结合，从而实现对船舶、海上平台等海洋结构物在强非线性波浪作用下的水动力性能进行综合分析。%

(3)	浮式平台波浪砰击、上浪等强非线性问题涉及复杂的气液固耦合现象，而本文并未考虑气相影响，%
并且对于砰击问题的准确预报需要使用更精确的物理模型。%
因此在未来工作中可以采用多相流SPH模型或其他多相流CFD方法进行耦合。
